\section{Introduction}
How small a container holds $n$ copies of a given shape? For unit squares in a square the question goes back to Erd\H{o}s and Graham \cite{erdosgraham}; catalogues now cover squares, triangles, hexagons and cubes in squares, triangles, circles and cubes \cite{friedmancenter,squaresproject}, and most entries are best known packings found by search rather than proved optimal. Most searches move rigid pieces: simulated annealing \cite{kirkpatrick}, perturbation and compression \cite{Gensane2004,gensane}, growing the pieces in size \cite{lubachevsky}, overlap minimisation in a shrinking container \cite{Umetani2009,Gardeyn2025}, or global solvers on exact non-overlap models \cite{Berthold2026}.

This paper changes the pieces instead. \emph{Hardening} starts every piece as the disk inscribed in it, compresses the disks under inward pressure on the container, and deforms them into the real shape while the pressure continues. It is graduated non-convexity \cite{blake} applied to geometry: a disk has no orientation, so the pieces are already densely packed when their corners appear and rotation becomes a choice. The shape family it moves along, a polygon dilated by a disk at fixed inradius, is not new: it is the spheropolygon of discrete-element modelling \cite{AlonsoMarroquin2008}, and Jin, Lu and Li moved dense periodic packings along its 3D analogue between tetrahedra and spheres \cite{Jin2015}. Jin, Lu and Li also ran it from spheres toward tetrahedra, in periodic packings. What we have not found is its use as an optimisation path toward a prescribed polygon inside a finite container, evaluated against tuned rigid searches for best known packings (Section~\ref{sec:related}). On squares it recovers several classical packings from random starts, and in 3D it found a packing of 12 unit cubes in a cube smaller than the best known one \cite{nakajima12}.

Here we ask whether that is a property of the path or of the problems we tried. We compare three paths through shape space (hardening, growing, and rigid pieces on the same schedule) and two classical baselines on seven families of shapes and containers (\totalruns{} runs in \ncampaigns{} campaigns), with matched budgets and comparable tuning effort. Since a packing search is judged by the best packing it ever finds, the comparison is by the \emph{lowest} container size each method reaches in a fixed number of runs, after constrained numerical tightening; the procedure being compared is the batch of runs, not a single run \cite{Birattari2005}. Our findings:
\begin{itemize}\itemsep1pt
\item No path is best everywhere. Untuned, hardening looked better than rigid starts; tuned with comparable effort, rigid starts on the same schedule reach the lowest size most often (Table~\ref{tab:main}).
\item Which path wins depends on the family. For squares in a square, hardening reaches the lowest size on \CsevenSquInSquHardenLowest{} of \CsevenSquInSquN{} instances and is the only method that low on \CsevenSquInSquHardenSole, including Trump's $n=11$ and Wainwright's $n=19$ packings; for triangles, rigid or growing starts are better (Section~\ref{sec:learned}).
\item Neither the disk-compressed start nor the area schedule alone reproduces this pattern (Table~\ref{tab:ablate}), and hardening's unique square packings end with more tilted pieces than rigid starts' best (Figure~\ref{fig:mechanism}). The pre-registered tests of these ablations are inconclusive on squares in a square.
\item On four families held out from all development, hardening is worse than rigid starts, and a pilot that chooses the path per instance is no better than always starting rigid (Table~\ref{tab:decision}).
\item An earlier indication that hardening gains more from longer runs did not replicate.
\end{itemize}
The result is a negative one with a specific exception, and the paper reports it as such. The study's other contribution is its record: every experiment was planned in writing before it ran, every run can be replayed, and every number in this paper is generated from an event log. Section~\ref{sec:matters} discusses where shape continuation could be useful beyond record tables.
